\documentclass[10pt, letterpaper]{article}
\usepackage{qrcode}
\usepackage{atveryend}
\usepackage{reddy_cover}
\usepackage{hyperref}

\begin{document}

\vspace{\headerSpacing}
\section*{COVER LETTER}

\vspace{\headerSpacing}
\begin{onecolentry}
	Aadarsha Gopala Reddy\\
	St. Louis, MO 63112\\
	(740) 802-1776\\
	adurs2002@gmail.com\\
\end{onecolentry}

\vspace{\headerSpacing}
\begin{onecolentry}
	July 21, 2026\\
\end{onecolentry}

\vspace{\headerSpacing}
\begin{onecolentry}
	Hiring Manager\\
	Associate Full-Stack Software Engineer\\
	Intramotev\\
	St. Louis, MO\\
\end{onecolentry}

\vspace{\headerSpacing}
\begin{onecolentry}
	\textbf{Subject:} Application for Associate Full-Stack Software Engineer - Intramotev
\end{onecolentry}

\vspace{\headerSpacing}
\begin{onecolentry}
	Dear Hiring Manager,
\end{onecolentry}

\vspace{\headerSpacing}
\begin{onecolentry}
	I am writing to apply for the Associate Full-Stack Software Engineer position at Intramotev. I completed my M.S. in Computer Science at Washington University in St.~Louis this past May, and I am based here in St.~Louis already. What drew me to this role specifically is the blend of mobile and embedded work behind it: building the operator-facing interface that runs on a physical tablet in the field, backed by services that talk to real vehicles, is exactly the kind of hardware-facing software I have gravitated toward on my own time.
\end{onecolentry}

\vspace{\headerSpacing}
\begin{onecolentry}
	As a Data Analyst Intern at Lab714, I designed and built the Figma-to-React UI for an all-in-one environmental sensor device, taking it from wireframes to a pixel-perfect, production-ready interface, and I built the AWS data ingestion pipelines that normalized telemetry from a deployed IoT device fleet and powered the real-time dashboards operators used to catch equipment alerts. More recently I built a real-time Android app in Kotlin that decodes deliberate eye blinks into Morse code using live camera-based facial-landmark tracking and a debounced state machine, and I wrote Arduino C++ firmware for an environment monitor that manages a hardware button interface, an LCD display, and a rolling EEPROM buffer that survives power cycles.
\end{onecolentry}

\vspace{\headerSpacing}
\begin{onecolentry}
	As an AI Engineer Intern at Crittero, I worked in TypeScript daily and built a configurable Python CLI test harness with deterministic seeding and CSV export so a TensorFlow/Keras recommendation model's evaluation would be reproducible across configurations. As a Graduate Teaching Assistant for a data engineering course this past semester, I taught fault-tolerant pipeline design and led students through Git-based collaborative workflows and code review. That testing discipline and Git workflow experience is what I would bring to building and testing C\# microservices connecting vehicles to operators.
\end{onecolentry}

\vspace{\headerSpacing}
\begin{onecolentry}
	I would welcome the chance to walk you through this work in more detail. My GitHub (\href{https://github.com/agopalareddy}{github.com/agopalareddy}) has the full source for the Android app and the embedded firmware project mentioned above. Thank you for considering my application; I look forward to the possibility of contributing to Intramotev's autonomous railcar fleet.
\end{onecolentry}

\vspace{\headerSpacing}
\begin{onecolentry}
	Sincerely,
\end{onecolentry}

\vspace{\headerSpacing}
\begin{onecolentry}
	Aadarsha Gopala Reddy\\
\end{onecolentry}
\end{document}
